\section{Native NLP Interior-Point Method: Theory, Algorithm, and Current Diagnostics}
\label{sec:engine-nlp-ipm}

This section documents the nonlinear interior-point implementation under
\texttt{include/hacdcpf/engine/ipm\_solver.hpp} and
\texttt{src/engine/ipm\_solver.cpp} with a narrower goal than
Section~\ref{sec:engine-solver-system}: to explain what the native NLP IPM is
mathematically solving, how the current code actually drives the iterations,
what the present test evidence says, and which observed behaviors should be
read as signs of a real algorithmic problem rather than mere tuning noise.

The discussion is intentionally diagnostic.  The current code already contains
multiple globalization paths, scaling hooks, restoration support, exact and
inexact curvature modes, and IPOPT fallback.  A reader who only sees the class
name \texttt{NativeIPMAdapter} can therefore easily overestimate its maturity.
The notebook below separates what is theoretically intended, what is already
implemented, and what the latest large-case evidence still says is fragile.

\subsection{Standard NLP Form and Barrier KKT System}

The canonical engine nonlinear program is the \texttt{NLPModel} from
\texttt{problem\_types.hpp}:
\begin{align}
  \min_x \quad & f(x) \\
  \text{s.t.}\quad & g(x) = 0, \\
                    & h(x) \le 0, \\
                    & \ell \le x \le u.
\end{align}

The engine converts all nonlinear inequalities and finite variable bounds into
one unified inequality block.  Introducing slacks
\(s > 0\) such that
\begin{equation}
  h(x) + s = 0,
\end{equation}
and multipliers \(\lambda\) for equalities and \(\mu > 0\) for inequalities,
the barrier KKT conditions are
\begin{align}
  \nabla f(x) + J_g(x)^\top \lambda + J_h(x)^\top \mu &= 0, \\
  g(x) &= 0, \\
  h(x) + s &= 0, \\
  S \mu &= \tau e,
\end{align}
where \(S = \operatorname{diag}(s)\).

The reduced Newton system used by both the parity-specific IPM and the generic
engine IPM is based on the primal Schur complement
\begin{equation}
  W = \nabla^2_{xx} \mathcal{L}(x,\lambda,\mu)
      + J_h(x)^\top \operatorname{diag}(\mu ./ s) J_h(x),
\end{equation}
with the Lagrangian
\begin{equation}
  \mathcal{L}(x,\lambda,\mu) = f(x) + \lambda^\top g(x) + \mu^\top h(x).
\end{equation}

This identity is critical when reading the current code quality:
if the solver only uses \(\nabla^2 f\) and drops the constraint-curvature term
inside \(\nabla^2_{xx} \mathcal{L}\), then the Newton direction is formally
incorrect for nonlinear constrained problems.  The recent engine changes fixed
that part by adding an explicit
\texttt{NLPModel::lagrangian\_hess(...)} callback and by using it inside the
native solver when available.

\subsection{Current Driver Architecture in \texttt{ipm\_solver.cpp}}

The file now contains two distinct native drivers behind the same adapter.

\paragraph{1. Filter driver (current default)}
The public \texttt{IPMOptions} now expose
\texttt{Globalization::Filter} as the default mode.  In this path, the solver
first optionally scales the problem, then runs a filter line-search method in
the style of W"achter--Biegler, with optional restoration and second-order
correction logic.  If the filter path fails, it can restore feasibility and
retry; if that still fails, it falls back to the Merit-backed native path and
ultimately to IPOPT.

The important consequence is practical rather than aesthetic:
``the native NLP IPM'' is not one single algorithm anymore.  It is a driver
stack with this order:
\begin{enumerate}
  \item scaled filter IPM,
  \item restoration-driven filter retry,
  \item merit-based native IPM,
  \item IPOPT fallback.
\end{enumerate}

\paragraph{2. Merit driver (fallback native path)}
The lower part of \texttt{ipm\_solver.cpp} still contains the earlier
Mehrotra-style merit-based driver.  It has recently been upgraded with:
\begin{itemize}
  \item exact Lagrangian Hessian support,
  \item residual-aware slack and multiplier initialization,
  \item normalized residual metrics,
  \item separate primal and dual step sizes, and
  \item best-iterate restoration.
\end{itemize}

These changes materially improved the behavior on large OPF-derived NLPs, but
they did not eliminate the main robustness issue: the generic merit line search
still breaks down on large parity OPF cases before the dedicated parity solver
does.

\subsection{Curvature Sources and Hessian Policy}

The current engine supports three curvature regimes.

\paragraph{Exact Lagrangian curvature}
If \texttt{lagrangian\_hess} is provided, the solver uses the exact
\(\nabla^2_{xx}\mathcal{L}\) callback directly.  This is now the preferred path
for parity OPF and for any NLP with nonlinear constraint curvature.

\paragraph{Objective-only Hessian}
If only \texttt{hess} is present, the solver uses the objective Hessian
\(\nabla^2 f\).  This is exact only when all constraints are linear; otherwise
it is a model misspecification, not just a performance compromise.

\paragraph{Recovered curvature}
If neither callback is present, the solver constructs sparse secant and
finite-difference approximations.  This path is still operational and is well
covered by the structured sparse-chain tests, but it is not the preferred path
for physics-heavy OPF NLPs.

\subsection{Algorithmic Workflow of the Merit Driver}

The current merit path in \texttt{ipm\_solver.cpp} can be summarized as follows.

\paragraph{Algorithm E5: Merit-backed Native NLP IPM (current implementation)}
\begin{algorithmic}[1]
\Procedure{\texttt{NativeIPM.solve\_nlp\_detail}}{$prob$}
  \State Validate callbacks, dimensions, and bounds
  \State Interiorize the initial point with respect to box bounds
  \State Evaluate \(f,\nabla f,g,J_g,h,J_h\)
  \State Initialize slacks and multipliers from inequality residuals
  \State Build exact or recovered Lagrangian Hessian
  \For{$k = 0,1,\dots$ until convergence or iteration limit}
    \State Form residuals \((r_{dual}, r_{eq}, r_{ineq})\)
    \State Compute normalized primal, dual, and complementarity metrics
    \State Build reduced primal system
    \State Solve affine predictor step
    \State Compute centering target and corrector step
    \State Compute separate primal and dual fraction-to-boundary steps
    \State Try merit-based backtracking on \((x,s,\lambda,\mu)\)
    \If{acceptable trial found}
      \State Accept step and continue
    \Else
      \State Recenter barrier state or terminate with line-search failure
    \EndIf
  \EndFor
  \If{best historical iterate satisfies tolerances}
    \State Restore best iterate and report convergence
  \Else
    \State Return failure or fall back to IPOPT
  \EndIf
\EndProcedure
\end{algorithmic}

The important diagnostic point is Step 10.  The generic solver still judges
trial steps with a scalar merit decrease condition.  For large OPF NLPs this is
currently weaker than the parity-specific direct primal/dual step acceptance,
and it is the main remaining failure mode after the recent scaling and
initialization fixes.

\subsection{Algorithmic Workflow of the Filter Driver}

The filter path is theoretically stronger and architecturally more ambitious.
It adds four ingredients that the earlier merit-only solver did not have:
\begin{itemize}
  \item problem scaling,
  \item filter-based globalization instead of one scalar merit test,
  \item optional second-order correction,
  \item one-shot restoration NLP construction and retry.
\end{itemize}

In theory this is the correct direction for difficult nonconvex NLPs.  In
practice the current source still has to be read as an implementation under
active hardening rather than a finished IPOPT-equivalent engine.  The notebook
should therefore be interpreted conservatively: the filter driver is the
intended production globalization path, while the merit driver remains the most
transparent reference path for code reading and regression triage.

\subsection{What the Unit Tests Currently Prove}

The current dedicated native-IPM test coverage is concentrated in
\texttt{tests/test\_nlp\_ipm\_solver.cpp} and
\texttt{tests/test\_minlp\_native\_ipm.cpp}.  The test evidence is summarized
in Table~\ref{tab:nlp-ipm-tests}.

\begin{table}[H]
\centering
\small
\caption{Native NLP IPM tests and what they actually verify}
\label{tab:nlp-ipm-tests}
\begin{tabularx}{\linewidth}{L{0.30\linewidth}L{0.22\linewidth}L{0.38\linewidth}}
\toprule
Test file / case & Status & What it proves \\
\midrule
\texttt{test\_nlp\_ipm\_solver}: small constrained quadratic & Passing & Basic correctness of the primal-dual solve path, multiplier output, and convergence on a smooth small NLP. \\
\texttt{test\_nlp\_ipm\_solver}: sparse chain NLP & Passing & Exact-Hessian and missing-Hessian paths both work on a structured sparse benchmark family. \\
\texttt{test\_nlp\_ipm\_solver}: exact Lagrangian Hessian callback & Passing & The solver now actually consumes \texttt{lagrangian\_hess} rather than silently dropping nonlinear constraint curvature. \\
\texttt{test\_minlp\_native\_ipm} & Passing & MINLP relaxations can call the native NLP IPM and return valid node relaxations for small cases. \\
Parity-specific Hessian finite-difference tests & Passing & The parity formulation's analytic Lagrangian Hessian matches finite-difference checks. \\
\bottomrule
\end{tabularx}
\end{table}

These tests are important but not sufficient.  They prove local algorithmic
correctness on small and structured problems; they do \emph{not} prove robust
globalization on large OPF NLPs.

\subsection{Structured Sparse Benchmark Results}

The synthetic benchmark family in
\texttt{benchmark/nlp\_ipm\_benchmark.cpp} still provides the strongest
evidence that the engine's linear algebra and Hessian handling are no longer
the main bottleneck.  On the structured sparse chain problems, the native exact
Hessian path and the missing-Hessian recovery path both solve successfully, and
the native solver remains competitive with or faster than IPOPT on this
benchmark family.

That result is still technically relevant: it means the core KKT machinery,
sparse Hessian handling, and callback contracts are not fundamentally broken.
If large OPF cases fail, the most likely culprit is now globalization and
problem scaling, not the absence of a Newton kernel.

\subsection{Large Parity-OPF Comparison Results}

The most informative recent evidence comes from the parity-OPF benchmark driver
in \texttt{tests/benchmark\_ac\_opf\_cases.cpp}, which now exposes three solve
modes on the same parity formulation:
\begin{itemize}
  \item \texttt{parity}: dedicated parity primal-dual IPM in
  \texttt{src/optimal_power_flow/parity\_ipm.cpp},
  \item \texttt{native}: generic \texttt{NativeIPMAdapter} applied to the same
  parity NLP callbacks,
  \item \texttt{ipopt}: IPOPT on the same parity NLP callbacks.
\end{itemize}

Table~\ref{tab:nlp-ipm-large-opf} summarizes the latest results after the
recent merit-driver upgrades.

\begin{table}[H]
\centering
\small
\caption{Large-case parity OPF comparison on the same NLP formulation}
\label{tab:nlp-ipm-large-opf}
\begin{tabularx}{\linewidth}{L{0.16\linewidth}L{0.13\linewidth}L{0.10\linewidth}L{0.15\linewidth}L{0.17\linewidth}L{0.19\linewidth}}
\toprule
Case & Mode & Converged & Mean time (ms) & Iterations & Status / interpretation \\
\midrule
IEEE118 & native & No & 2307.323 & 90 & \texttt{NativeIPM line search failed}; much better than the earlier 500-iteration stall, but still not robust. \\
IEEE118 & ipopt & No & 1475.001 & 500 & \texttt{Max iterations exceeded}; remains near the correct objective and near-feasible. \\
IEEE118 & parity & Yes & 15880.952 & 23 & Dedicated parity solver converges and passes PF cross-validation. \\
PEGASE2869 & native & No & 76498.973 & 33 & \texttt{NativeIPM line search failed}; materially improved over the earlier 500-iteration collapse, but still unusable as a final OPF solver. \\
PEGASE2869 & ipopt & No & 65131.865 & 500 & \texttt{Max iterations exceeded}; stays close to the parity objective. \\
PEGASE2869 & parity & Yes & 123052.538 & 61 & Dedicated parity solver converges and passes PF cross-validation. \\
\bottomrule
\end{tabularx}
\end{table}

Two conclusions follow immediately.
\begin{enumerate}
  \item The recent engine changes are real improvements: the generic native NLP
  IPM no longer simply burns the full iteration budget with catastrophic
  residuals on these OPF cases.
  \item The generic native NLP IPM is still not production-robust on large OPF
  NLPs.  The dominant failure is now a globalization failure, not missing
  curvature.
\end{enumerate}

\subsection{Current Diagnostic Interpretation: Is the Code Still Wrong?}

For code review purposes, the current answer is nuanced.

\paragraph{What is already correct enough to trust}
\begin{itemize}
  \item The engine now has a mathematically correct interface for exact
  Lagrangian Hessians.
  \item Small and structured NLP tests pass on both exact and recovered
  curvature paths.
  \item The parity formulation's Hessian and Jacobian contracts have dedicated
  finite-difference validation.
  \item The generic solver's latest upgrades improved large OPF behavior by a
  large margin compared with the earlier 500-iteration stalls.
\end{itemize}

\paragraph{What still indicates a real algorithmic problem}
The remaining large-case failures are not benign.  The following behaviors
should be read as evidence that the generic NLP IPM is still incomplete for
hard OPF-class problems:
\begin{itemize}
  \item repeated \texttt{line search failed} termination on large parity NLPs,
  \item native objective values far away from the parity and IPOPT trajectories,
  \item generic native failure on cases where the dedicated parity solver
  converges on the same formulation,
  \item dependence on fallback or parity-specific kernels for final robustness.
\end{itemize}

In other words, the code is no longer ``wrong'' in the narrow sense of missing
constraint curvature or having no native nonlinear interior-point core.
However, it is still incomplete in the stronger sense required for large-scale
OPF robustness.  The current bottleneck is algorithmic globalization rather
than Hessian formation.

\subsection{Most Likely Remaining Gap}

The recent parity-derived changes fixed the easiest-to-identify deficiencies:
better barrier-state initialization, normalized residuals, exact Lagrangian
curvature, and separate primal/dual step sizes.  The strongest remaining gap is
that the generic merit path still uses a scalar merit backtracking test, while
the dedicated parity solver accepts fraction-to-boundary primal and dual steps
more directly and augments them with OPF-specific corrector logic.

So the most likely next successful improvement path is:
\begin{enumerate}
  \item harden the filter-based driver until it reliably dominates the merit
  fallback on large OPF cases,
  \item or, alternatively, import the parity solver's direct primal/dual step
  acceptance and extra-corrector behavior more faithfully into the generic
  native path.
\end{enumerate}

For a reviewer trying to judge the present source tree, the practical summary
is simple: the native NLP IPM is now a real solver and not a placeholder, but
its robustness envelope is still substantially narrower than the dedicated
parity OPF IPM on large network NLPs.

\subsection{Reviewer-Style Risk Checklist}

For code review and regression triage, the current source can be reduced to the
following seven risk points.  These are the items most likely to explain why
the native NLP IPM still underperforms on large OPF-class problems.

\begin{enumerate}
  \item \textbf{Globalization mismatch remains the primary risk.}
  The generic native path still relies on merit or filter globalization logic
  that is materially more fragile than the direct fraction-to-boundary step
  acceptance used by the dedicated parity IPM.  Current large-case failures are
  still dominated by \texttt{line search failed} or earlier filter-stage
  rejection, not by Hessian unavailability.

  \item \textbf{Filter driver maturity is still below production level for OPF.}
  The default path in \texttt{ipm\_solver.cpp} is now the filter-based driver,
  but the current large-case evidence does not yet show that it dominates the
  merit fallback or IPOPT.  A code reviewer should therefore treat the filter
  driver as an active development path rather than as a solved final design.

  \item \textbf{The merit fallback is improved but still not robust enough.}
  The recent parity-derived upgrades fixed barrier initialization, residual
  normalization, exact Lagrangian curvature usage, separate primal and dual step
  lengths, and best-iterate restoration.  Even so, the merit path still fails
  on IEEE118 and PEGASE2869 when applied to the parity NLP callbacks.

  \item \textbf{Performance is still problem-class dependent rather than broadly solved.}
  The sparse structured chain benchmarks show that the linear algebra kernel and
  Hessian plumbing are not fundamentally broken.  However, those wins do not yet
  transfer to large nonconvex OPF NLPs.  A reviewer should therefore reject any
  blanket claim that the native NLP IPM is already competitive across realistic
  nonlinear planning workloads.

  \item \textbf{The generic solver still trails the dedicated parity solver on the same formulation.}
  This is the strongest practical red flag in the current tree.  When the same
  parity NLP callbacks are solved by \texttt{parity}, \texttt{native}, and
  \texttt{ipopt}, the dedicated parity solver converges while the generic native
  solver still fails.  That means the remaining gap is algorithmic and not just
  an interface or model-building issue.

  \item \textbf{Fallback layering can hide native weakness.}
  The current adapter stack can flow through filter solve, restoration retry,
  merit solve, and finally IPOPT fallback.  This is a useful resilience feature,
  but it also means that a superficially successful solve may not actually prove
  native robustness unless the solver path is checked explicitly.

  \item \textbf{The present bottleneck is not ``missing Hessian'' anymore.}
  That diagnosis was accurate earlier, but it is no longer the main one.  The
  more serious current risks are globalization quality, scaling fidelity, and
  parity-to-generic algorithm mismatch in difficult OPF geometries.
\end{enumerate}

In short, the present code base should be read as follows: the native NLP IPM
is algorithmically real and test-backed, but its performance and robustness are
still clearly insufficient for large OPF-class NLPs.  This is now a
globalization-and-driver-design problem more than a second-derivative problem.